\section{An exact theorem over every finite coefficient field}
\label{sec:scalar}

\subsection{Setup and statement}
Fix a prime $p$, a finite extension $K/\Q_p$, and the valuation $v$ normalized by $v(p)=1$ on $\C_p$. Write
\[
e=e(K/\Q_p),\quad w=v_p(e),\quad
N_h=(p-1)p^{h-1},\quad r_h=N_h^{-1}\quad(h\ge1).
\]
Let $\mathcal Z_h=\{\zeta-1:\zeta\text{ has exact order }p^h\}$.
Every point of $\mathcal Z_h$ has valuation $r_h$.

For a nonzero polynomial $g=\sum_jc_jT^j$ and $\rho\ge0$, its Gauss valuation is
\[
G_\rho(g)=\min_j\{v(c_j)+j\rho\}.
\]
Only nonzero coefficients are included in this minimum. The Gauss valuation is multiplicative. If $v(x)\ge\rho$, then $v(g(x))\ge G_\rho(g)$.

\begin{definition}
A packet consists of finitely many nonzero $g_i\in\OO_K[T]$ and finite thresholds $t_i\ge G_\rho(g_i)$. Its budget is $\sum_i t_i$. At $z$ define
\[
D(z)=\sum_i\max\{0,v(g_i(z))-t_i\}.
\]
A level $h$ is blocked at zero tolerance if $D(z)>0$ for every $z\in\mathcal Z_h$.
\end{definition}
Actual witnesses $x_i$ with $v(x_i)\ge\rho$ and $t_i=v(g_i(x_i))<\infty$ give a special case. Allowing general Gauss thresholds is convenient for the multivariable induction.

\begin{theorem}[Exact zero-tolerance obstruction cost]
\label{thm:scalar}
Let $S$ be a nonempty finite subset of the positive integers, and suppose
$\rho\ge\max_{h\in S}r_h$. Among packets over $\OO_K$ that block every $h\in S$, the minimum budget is
\[
 C_K(S)=\sum_{h\in S}a_h,\qquad
 a_h=p^{-\min(w,h-1)}.
\]
The minimum is attained by one explicit binomial per level with common witness $0$. Consequently, if $\sum_i t_i\le B<C_K(S)$, there are $h\in S$ and $z\in\mathcal Z_h$ for which
\[
v(g_i(z))\le t_i\quad\text{for every }i.
\]
The strict crossing $B<C_K(S)$ is necessary for a universal guarantee.
\end{theorem}

The exact answer depends on the wild part $p^w$ of $e$. In particular, every tamely ramified coefficient field has the same zero-tolerance obstruction cost as $\Q_p$. This assertion concerns zero tolerance; positive tolerances can distinguish finer properties of the field.

\subsection{The root budget across distinct levels}
\begin{lemma}[A root cannot spend its budget twice]
\label{lem:root-budget}
Let $M_h(g)$ count the roots of $g$ having valuation $r_h$, with multiplicity. If $g\in\OO_K[T]\setminus\{0\}$ and $\rho\ge\max_{h\in S}r_h$, then
\[
\sum_{h\in S}r_hM_h(g)\le G_\rho(g).
\]
\end{lemma}
\begin{proof}
Factor $g=A\prod_\alpha(T-\alpha)$ over the algebraic closure, counting multiplicities and using $v(0)=+\infty$. Multiplicativity gives
\[
 G_\rho(g)-G_0(g)
 =\sum_\alpha\bigl(\min\{\rho,v(\alpha)\}-\min\{0,v(\alpha)\}\bigr).
\]
Every summand is nonnegative. A root of valuation $r_h$ contributes exactly $r_h$, and distinct $h$ have distinct radii. Finally $G_0(g)\ge0$ by integrality.
\end{proof}

For completeness, multiplicativity does not require guessing a generic evaluation point: in a product, choose the smallest exponents attaining the weighted coefficient minima in the two factors. Their sum is the smallest possible exponent of minimum weight in the product, and its minimum-weight contribution cannot cancel. This proves the usual Gauss-valuation identity.

\subsection{Ramification forces residual multiplicity}
\begin{lemma}[Multiplicity at one torsion radius]
\label{lem:initial}
Fix $h$ and $z_0\in\mathcal Z_h$. For $g=\sum c_jT^j\ne0$, choose $b$ in a finite extension with $v(b)=G_{r_h}(g)$ and reduce
\[
\mathcal I_g(Y)=\overline{b^{-1}g(z_0Y)}
\]
in the residue field of a common finite extension. Then:
\begin{enumerate}
\item $\mathcal I_g$ is nonzero and its number of nonzero roots, with multiplicity, is $M_h(g)$.
\item Every nonzero root has multiplicity at least
\[
 P_h=p^{\max(h-1-w,0)}.
\]
\item If $z\in\mathcal Z_h$ and $v(g(z))>t\ge G_\rho(g)$, with $\rho\ge r_h$, then $\overline{z/z_0}$ is a root of $\mathcal I_g$.
\end{enumerate}
\end{lemma}
\begin{proof}
A coefficient survives reduction exactly when
$v(c_j)+j/N_h=G_{r_h}(g)$. If both indices $j,k$ survive, then
\[
 \frac{j-k}{N_h}\in e^{-1}\Z,
 \qquad \frac{N_h}{\gcd(e,N_h)}\mid j-k.
\]
The $p$-part of the last divisor is $P_h$. After removing a monomial, $\mathcal I_g$ is therefore a polynomial in $Y^{P_h}$. Over an algebraic closure of the residue field, characteristic $p$ makes it a $P_h$-th power. Every nonzero root has the asserted multiplicity.

The first assertion also follows directly from the factorization of $g$. After scaling at radius $r_h$, a root of smaller valuation contributes a nonzero constant to the initial polynomial, a root of larger valuation contributes a monomial $Y$, and a root of exactly that valuation contributes a nonzero linear root. Thus the difference between the largest and smallest surviving exponents, equivalently the nonzero-root degree, is precisely $M_h(g)$.

For the last assertion, monotonicity gives $t\ge G_\rho(g)\ge G_{r_h}(g)$. Strict excess makes $b^{-1}g(z)$ reduce to zero. The normalized coordinate $z/z_0$ is a unit, so its residue is a nonzero root of $\mathcal I_g$.
\end{proof}

\subsection{Proof of the lower bound}
Write $z_0=\zeta_0-1$ with $\zeta_0$ primitive of order $p^h$. For $a=1,\ldots,p-1$ the points
\[
z_a=\zeta_0^a-1
\]
are primitive and
\[
 \overline{z_a/z_0}
 =\overline{1+\zeta_0+\cdots+\zeta_0^{a-1}}
 =a\in\F_p^\times.
\]
If level $h$ is blocked, each of these $p-1$ residual directions is a root of at least one packet member's initial polynomial. By Lemma~\ref{lem:initial},
\[
 \sum_i M_h(g_i)\ge(p-1)P_h.
\]
Multiplying by $r_h$ gives
\[
 r_h\sum_i M_h(g_i)
 \ge\frac{(p-1)p^{\max(h-1-w,0)}}{(p-1)p^{h-1}}
 =p^{-\min(w,h-1)}=a_h.
\]
Summing over $S$ and applying Lemma~\ref{lem:root-budget} proves
\[
 \sum_{h\in S}a_h
 \le\sum_i\sum_{h\in S}r_hM_h(g_i)
 \le\sum_iG_\rho(g_i)
 \le\sum_i t_i.
\]
This establishes the claimed lower bound without assuming that the packet's roots themselves are roots of unity.

\subsection{Sharp binomial obstructions}
Fix a level $h$ and set
\[
 k=\min(w,h-1),\qquad d_h=N_h/p^k.
\]
Choose a uniformizer $\pi$ of $K$. Since $p^k\mid e$, an element
$c_h=u\pi^{e/p^k}$ has valuation $p^{-k}$. The residue field of $K$ is perfect, so one can choose a unit $u$ for which
\[
 \overline{c_h^{p^k}/(-p)}=1.
\]
Explicitly, choose its residue to solve
$\overline u^{\,p^k}=\overline{(-p)/\pi^e}$; Frobenius is bijective on the finite residue field.

Now put
\[
g_h(T)=T^{d_h}-c_h,\qquad t_h=v(g_h(0))=p^{-k}.
\]
The shifted cyclotomic polynomial satisfies
\[
 \Phi_{p^h}(1+T)\equiv T^{N_h}\pmod p,
 \qquad \Phi_{p^h}(1)=p.
\]
All its intermediate terms have valuation strictly larger than $1$ when evaluated at $z\in\mathcal Z_h$. Hence
\[
 \overline{z^{N_h}/(-p)}=1.
\]
It follows that the unit $z^{d_h}/c_h$ has $p^k$-th power with residue $1$. In characteristic $p$ the only $p^k$-th root of $1$ is $1$, so
\[
v(g_h(z))>v(c_h)=t_h
\quad(z\in\mathcal Z_h).
\]
The Gauss baseline is also $G_\rho(g_h)=t_h$, since $\rho\ge r_h$. At any other level $h'$ the two terms in $g_h(z)$ have different valuations, and
\[
v(g_h(z))=\min(d_hr_{h'},t_h)\le t_h.
\]
Thus this binomial blocks exactly its designated level with budget $a_h$. Taking one for every $h\in S$ attains $C_K(S)$ and completes the proof of Theorem~\ref{thm:scalar}.

\subsection{Optimal cutoffs and a density version}
\begin{corollary}[Optimal consecutive-level endpoint]
\label{cor:cutoff}
For $S=[h_0,H]\cap\Z$, the least universal endpoint is the first $H\ge h_0$ satisfying
\[
 \sum_{h=h_0}^{H}p^{-\min(w,h-1)}>B.
\]
In particular, when $h_0\ge w+1$ it is exactly
\[
 H=h_0+\floor{p^wB}.
\]
For every $h_0$, this last expression remains a sufficient endpoint.
\end{corollary}
Equality of cost with $B$ still permits the sharp obstruction. The supplied code uses exact rational arithmetic for this strict crossing.

\begin{corollary}[A weighted sampling guarantee]
\label{cor:density}
Choose a finite set $S$ as in Theorem~\ref{thm:scalar}, choose $h$ with probability $a_h/C_K(S)$, then choose $z$ uniformly from $\mathcal Z_h$. For any packet of budget at most $B$, the probability of satisfying every separate constraint is at least
\[
 \max\{0,1-B/C_K(S)\}.
\]
\end{corollary}
\begin{proof}
Put $b_h=r_h\sum_iM_h(g_i)$. At most $\sum_iM_h(g_i)/P_h$ of the $p-1$ residual directions can be roots of any initial polynomial. Every direction contains $p^{h-1}$ primitive points. The proportion of failing points at level $h$ is at most
\[
 \min\{1,b_h/a_h\}.
\]
The weighted failure probability is therefore at most
$\sum_{h\in S}b_h/C_K(S)\le B/C_K(S)$.
\end{proof}
This is a probability bound on a finite search. It does not assert a bit-complexity bound for evaluating arbitrary algebraic input polynomials.

\begin{center}
\small
\begin{tabular}{rrrrrrr}
\toprule
$p$ & $e$ & $s=v_p(e)$ & $h_0$ & $B$ & Old $H$ & Exact $H$\\
\midrule
2 & 1 & 0 & 1 & 3 & 4 & 4\\
2 & 105 & 0 & 1 & 3 & 316 & 4\\
3 & 100 & 0 & 1 & 3 & 301 & 4\\
5 & 1008 & 0 & 1 & 3 & 3025 & 4\\
2 & 8 & 3 & 1 & 3 & 25 & 14\\
2 & 24 & 3 & 1 & 3 & 73 & 14\\
3 & 27 & 3 & 1 & 2 & 55 & 19\\
3 & 54 & 3 & 1 & 2 & 109 & 19\\
5 & 100 & 2 & 1 & 3 & 301 & 48\\
2 & 12 & 2 & 3 & 1 & 15 & 7\\
2 & 8 & 3 & 5 & 3 & 29 & 29\\
3 & 100 & 0 & 1 & 1/2 & 51 & 1\\
\bottomrule
\end{tabular}

\end{center}
Here the old cutoff is $h_0+\lfloor eB\rfloor$ from the previous fixed-field argument. The table's exact cutoff uses the full transient sum, not only the constant-tail estimate.

\subsection{Relationship to earlier work}
\label{sec:scalar-literature}
The unrestricted-coefficient minimax formula and the exact cyclotomic-coefficient formula were already proved in the supplied reports. The latter specializes at zero tolerance to our answer when $K=\Q_p(\zeta_{p^s})$: then $w=s-1$ and the shell cost is $p^{1-\min(h,s)}$. Thus the new formula agrees on the overlap and extends that zero-tolerance answer to arbitrary finite fields.

The proof uses the same root-budget idea as the upstream torsion selector, but replaces unrestricted nearest-root counting by multiplicity forced by the coefficient value group. Bhargava's $P$-ordering and fixed-divisor framework provides relevant classical context \cite{bhargava,fixdiv}. Our assertion concerns a particular simultaneous positive-excess obstruction, including the distribution of one total budget across different cyclotomic radii. We have not established that this exact formulation is absent from all prior literature.

The positive-tolerance optimum over a general noncyclotomic coefficient field remains open in this report. Zero tolerance is precisely the regime needed for the exact individual specialization theorem below.

